%% ============================================================
%%  Chapter 4 - Formal Mathematical Model
%% ============================================================
\chapter{Formal Mathematical Model}
\label{ch:math}

\section{System Variables}
\label{sec:sysvars}

Let $\mathcal{S}=\{s_1,\ldots,s_N\}$ denote subscribers and
$\mathcal{A}=\{a_1,\ldots,a_M\}$ participating access points. Let
$\mathcal{K}$ denote traffic slices, $\mathcal{D}$ administrative domains,
and $\mathcal{H}$ available backhaul modes. Subscriber $s_i$ is represented by
\begin{equation}
  s_i=(id_i,d_i,p_i,b_i^{min},b_i^{max},\ell_i^{max},\delta_i^{max},\alpha_i,\rho_i),
  \label{eq:subscriber-vector}
\end{equation}
where $id_i$ is pseudonymised identity, $d_i$ home domain, $p_i$ policy profile,
$b_i^{min},b_i^{max}$ bandwidth bounds, $\ell_i^{max}$ latency bound,
$\delta_i^{max}$ loss bound, $\alpha_i$ service class, and $\rho_i$ accounting
identifier.

AP $a_j$ is represented by
\begin{equation}
  a_j=(cap_j,bh_j,c_j,\eta_j,\phi_j,\omega_j,\Redge_j),
  \label{eq:ap-vector}
\end{equation}
where $cap_j$ is wireless capacity, $bh_j$ effective backhaul capacity,
$c_j$ channel, $\eta_j$ utilisation, $\phi_j$ local policy, $\omega_j$
survivability state, and $\Redge_j$ edge resource vector.

The assignment variables are
\begin{equation}
  x_{ij}=
  \begin{cases}
    1, & s_i\text{ is served by }a_j,\\
    0, & \text{otherwise,}
  \end{cases}
  \qquad
  z_{ik}=
  \begin{cases}
    1, & s_i\text{ is mapped to slice }k,\\
    0, & \text{otherwise.}
  \end{cases}
  \label{eq:assignment-vars}
\end{equation}

A subscriber session is valid only if
\begin{equation}
  Valid(o_i)=1\Longleftrightarrow Auth_i=1\land Pol_i=1\land Iso_i=1\land Acct_i=1.
  \label{eq:valid-session}
\end{equation}

\subsection{Notation, Parameter Selection, and Normalisation Convention}
\label{sec:notation-parameter-policy}

To avoid ambiguity between modelling variables, calibrated coefficients, and
experimentally fixed thresholds, this section defines the notation and
parameter-selection conventions used throughout the mathematical model and evaluation
framework. Parameters are classified into four groups: operational policy parameters,
measured quantities, calibrated coefficients, and pre-registered evaluation thresholds.

\begin{table}[h]
\centering
\caption{Additional notation and parameter-selection conventions.}
\label{tab:parameter_conventions}
\begin{tabular}{p{0.18\linewidth} p{0.52\linewidth} p{0.22\linewidth}}
\hline
Symbol & Meaning & Selection method \\
\hline
$T_{\mathrm{rot}}$ &
Nominal rotation period of the cached-token signing key $K_{\mathrm{net}}$. &
Operational security policy. \\

$T_{\mathrm{token}}$ &
Maximum validity window of a cached offline-access token. &
Chosen such that $T_{\mathrm{token}} < T_{\mathrm{rot}}$. \\

$\Delta_{\min}$ &
Minimum practically meaningful improvement in the composite evaluation score. &
Fixed before comparative experiments. \\

$B_{\mathrm{curr}}$ &
Current measured available backhaul capacity during a telemetry epoch. &
Measured from AP/gateway telemetry. \\

$B_{\mathrm{nominal}}$ &
Nominal expected backhaul capacity for the AP or scenario class. &
Configured baseline or pilot-free engineering value. \\

$\kappa_B$ &
Sensitivity of telemetry interval to backhaul degradation. &
Calibrated coefficient. \\

$\kappa_{\mathrm{CPU}}$ &
Sensitivity of telemetry interval to CPU pressure. &
Calibrated coefficient. \\

$\kappa_{\mathrm{RTT}}$ &
Sensitivity of the LEO macro-epoch backoff to RTT variance. &
Calibrated coefficient. \\

$\kappa_{\mathrm{drop}}$ &
Sensitivity of the LEO macro-epoch backoff to packet-drop rate. &
Calibrated coefficient. \\
\hline
\end{tabular}
\end{table}

The cached-token key rotation parameter $T_{\mathrm{rot}}$ is treated as an
operational security-policy parameter rather than a queueing or optimisation
parameter. Cached-token validity is bounded by a shorter token lifetime
$T_{\mathrm{token}}$, with
\begin{equation}
  T_{\mathrm{token}} < T_{\mathrm{rot}}.
  \label{eq:token-rotation-bound}
\end{equation}
During local degraded operation, cached-token authentication is rejected when the
locally held key age exceeds $2T_{\mathrm{rot}}$. This conservative bound prevents
prolonged use of stale token-signing material while allowing short temporary service
continuity during backhaul or AAA unavailability.

The threshold $\Delta_{\min}$ in the acceptance criterion is fixed before comparative
baseline experiments are performed. It represents the minimum practically meaningful
improvement in the normalised composite score $S(M)$, not a value tuned after
observing results. Unless otherwise justified by pilot-free engineering constraints, the
default value is
\begin{equation}
  \Delta_{\min}=0.05.
  \label{eq:delta-min-default}
\end{equation}
Sensitivity analysis will report whether the main conclusions change for
\begin{equation}
  \Delta_{\min}\in\{0.02,0.05,0.10\}.
  \label{eq:delta-min-sensitivity}
\end{equation}

All convex weighted sums use non-negative weights normalised to one. Specifically,
the graph-interference weights $\beta_r$, CMDP cost weights $\omega_r$, LEO cost
weights $\psi_r$, and composite-score weights $\zeta_r$ satisfy
\begin{align}
  &\beta_r\geq0,\qquad \sum_r\beta_r=1, \label{eq:beta-normalisation}\\
  &\omega_r\geq0,\qquad \sum_r\omega_r=1, \label{eq:omega-normalisation}\\
  &\psi_r\geq0,\qquad \sum_r\psi_r=1, \label{eq:psi-normalisation}\\
  &\zeta_r\geq0,\qquad \sum_r\zeta_r=1. \label{eq:zeta-normalisation}
\end{align}
These weights are fixed within a declared scenario class before comparative
evaluation. They are not claimed to be globally transferable across dense urban,
semi-urban, rural-extension, and mountain-backhaul-constrained scenarios without
recalibration.

The telemetry-throttle and LEO-backoff coefficients are treated differently from
convex cost weights. The coefficients $\kappa_B$, $\kappa_{\mathrm{CPU}}$,
$\kappa_{\mathrm{RTT}}$, and $\kappa_{\mathrm{drop}}$ are non-negative sensitivity
parameters:
\begin{equation}
  \kappa_B\geq0,\qquad
  \kappa_{\mathrm{CPU}}\geq0,\qquad
  \kappa_{\mathrm{RTT}}\geq0,\qquad
  \kappa_{\mathrm{drop}}\geq0.
  \label{eq:kappa-nonnegative}
\end{equation}
They control how strongly telemetry sampling and LEO macro-epoch backoff react to
backhaul degradation, CPU pressure, RTT variance, and packet-drop rate,
respectively. They are calibrated through controlled sensitivity analysis and are not
required to sum to one.

\section{Finite Implementation Queue and Auxiliary LDT Queue}
\label{sec:queue-split}

The model separates the physical finite-buffer implementation queue from the auxiliary
infinite-buffer queue used for large-deviation analysis. This avoids the contradiction of
imposing a finite upper boundary $B_k$ while later evaluating $P(Q_k>B_k)$.

\subsection{Finite Drop-Tail Implementation Model}
\label{sec:finite-queue}

For slice $k$, the implemented queue is a reflected/censored workload process
$Q_k(t)\in[0,B_k]$. The service rate used in the queueing model is not the nominal
configured traffic-control rate alone. Instead, the effective bottleneck service capacity
is
\begin{equation}
  C_k^{eff}(t)=
  \min\{C_k^{tc}(t),C_k^{air}(t),C_k^{cpu}(t),C_k^{tun}(t),C_k^{bh}(t)\},
  \label{eq:effective-service}
\end{equation}
where $C_k^{tc}(t)$ is the configured traffic-control service rate,
$C_k^{air}(t)$ is the airtime-limited wireless capacity, $C_k^{cpu}(t)$ is the
CPU-limited forwarding capacity, $C_k^{tun}(t)$ is the tunnel-limited capacity after
encapsulation or encryption overhead, and $C_k^{bh}(t)$ is the available backhaul
capacity.

The finite-buffer workload process is therefore modelled as
\begin{equation}
  dQ_k(t)=
  \left(\lambda_k(t)-C_k^{eff}(t)\right)dt
  +\sigma_k dW_t+dR_k^0(t)-dL_k^B(t),
  \label{eq:finite-queue-sde}
\end{equation}
where $R_k^0(t)$ is the lower-bound reflection process at zero and $L_k^B(t)$ is the
upper-bound loss/local-time process at buffer threshold $B_k$. Overflow in the
implementation model is measured through the loss process $L_k^B(t)$, not through
the impossible event $Q_k(t)>B_k$.

The associated probability-flux expression for the unconstrained interior is
\begin{equation}
  J(q,t)=\mu_k(t)p(q,t)-\frac{\sigma_k^2}{2}\frac{\partial p(q,t)}{\partial q},
  \qquad
  \mu_k(t)=\lambda_k(t)-C_k^{eff}(t).
  \label{eq:flux}
\end{equation}
At $q=0$, reflection imposes $J(0,t)=0$. At $q=B_k$, finite-buffer packet loss is
recorded through the loss process rather than by claiming that the physical queue
exceeds $B_k$. All queueing, effective-bandwidth, and SLA calculations use
$C_k^{eff}(t)$, not the nominal configured traffic-control rate.

\subsection{Diffusion Coefficient Estimation}
\label{sec:diffusion}

The diffusion coefficient must be a variance rate, not a raw telemetry-window variance.
Let $A_k(t)$ denote cumulative input to slice $k$. For macro epoch $\tau$, define the
quasi-static interval
\begin{equation}
  T_\tau=[\tau T_{GNN},(\tau+1)T_{GNN}).
  \label{eq:macro-epoch}
\end{equation}
Within $T_\tau$, the empirical diffusion estimate is
\begin{equation}
  \hat{\sigma}_{k,\tau}^2
  =
  \frac{1}{\Delta}
  \operatorname{Var}
  \left[
    A_k(t+\Delta)-A_k(t)-\hat{\lambda}_{k,\tau}\Delta
  \right],
  \qquad t,t+\Delta\in T_\tau,
  \label{eq:diffusion-estimator}
\end{equation}
where $\Delta$ is the measurement interval and $\hat{\lambda}_{k,\tau}$ is the
estimated mean input rate for slice $k$ during epoch $\tau$. This definition gives
units of buffer-units squared per second and is compatible with the drift-diffusion
approximation.

The diffusion coefficient $\sigma_k$ in Eq.~\eqref{eq:finite-queue-sde} is interpreted
at the macro-epoch scale. In implementation and empirical estimation, $\sigma_k^2$ is
instantiated by the epoch-indexed estimator $\hat{\sigma}^{2}_{k,\tau}$ in
Eq.~\eqref{eq:diffusion-estimator} under the quasi-static assumption that arrival and
effective-service statistics remain approximately stationary within $T_\tau$.

\section{Large Deviations Bound for the Auxiliary Infinite Queue}
\label{sec:ldt}

Large Deviations Theory is applied to an auxiliary infinite-buffer workload
$Q_k^{\infty}(t)$, not to the finite implementation queue. The virtual tail probability
\begin{equation}
  P(Q_k^{\infty}>B_k)
  \label{eq:virtual-tail-probability}
\end{equation}
approximates the overflow risk that the finite queue would experience at threshold
$B_k$.

Let $A_k(t)$ be the cumulative input process. Its effective bandwidth is
\begin{equation}
  E_k(\theta)
  =
  \lim_{t\to\infty}
  \frac{1}{\theta t}
  \log\mathbb{E}\left[e^{\theta A_k(t)}\right],
  \qquad \theta>0.
  \label{eq:eff-bw-def}
\end{equation}
The stability condition is $E_k(0^+)<C_k^{eff}$, where $E_k(0^+)$ is the mean input
rate when the limit exists. The decay exponent is defined safely as
\begin{equation}
  \theta_k^*
  =
  \sup\{\theta>0:E_k(\theta)<C_k^{eff}\}.
  \label{eq:sup-root}
\end{equation}
If $E_k(\theta)$ is continuous, strictly increasing, and crosses $C_k^{eff}$, then the
supremum is the unique positive solution of
\begin{equation}
  E_k(\theta_k^*)=C_k^{eff}.
  \label{eq:root-conditional}
\end{equation}
Otherwise, the analysis reports the supremum-based bound rather than assuming root
existence.

Under the effective-bandwidth assumptions, the auxiliary workload satisfies the
logarithmic asymptotic relation
\begin{equation}
  \lim_{B_k\to\infty}
  \frac{1}{B_k}
  \log P(Q_k^{\infty}>B_k)
  =
  -\theta_k^*.
  \label{eq:ldt-asymptotic}
\end{equation}
For finite $B_k$, the quantity $e^{-\theta_k^*B_k}$ is used only as an
exponential-order overflow-risk indicator, not as a calibrated finite-buffer loss
probability unless validated against measured queue-loss traces.

This bound estimates overflow risk only. Cross-slice isolation is not a consequence of
large-deviation theory; it is enforced by VLAN/VNI separation, queue-discipline
configuration, and verified packet tracing in the edge datapath.

\section{Hard Safety Feasibility and State-Dependent Safe Actions}
\label{sec:pisafe}

\begin{definition}[State-dependent safe action set]
For a state $x\in\mathcal{X}$, the safe action set is
\begin{equation}
  \Usafe(x)=
  \{u\in\mathcal{U}(x):
  Auth(x,u)=1,\;
  Pol(x,u)=1,\;
  Iso(x,u)=1,\;
  Acct(x,u)=1,\;
  EdgeFeas(x,u)=1\}.
  \label{eq:usafe-state}
\end{equation}
\end{definition}

The edge-feasibility predicate is
\begin{equation}
\begin{aligned}
  EdgeFeas(x,u)=1
  \Longleftrightarrow\;&
  CPU_j\leq CPU_j^{max}
  \land RAM_j\leq RAM_j^{max}
  \land IRQ_j\leq IRQ_j^{max} \\
  &\land \Omega_j\leq\Omega_j^{max}
  \land C_k^{eff}>E_k(0^+).
\end{aligned}
\label{eq:edge-feasibility}
\end{equation}
The final condition ensures that the effective service capacity of the selected slice
exceeds the mean input rate implied by the effective-bandwidth model. Thus, a safe
action is not merely authenticated and policy-valid; it must also be feasible under
measured edge-resource and queue-stability constraints.

\begin{definition}[Safe policy set]
A policy is safe if it chooses only safe actions in every reachable state:
\begin{equation}
  \Pisafe=
  \{\pi:\pi(x)\in\Usafe(x),\;\forall x\in\mathcal{X}\}.
  \label{eq:pisafe-corrected}
\end{equation}
\end{definition}

Basic access constraints are
\begin{align}
  &\sum_{j=1}^M x_{ij}\leq1, &&\forall i,
  \label{eq:single-ap}\\
  &x_{ij}\leq Auth_i, &&\forall i,j,
  \label{eq:auth-req}\\
  &z_{ik}\leq Pol_{ik}, &&\forall i,k.
  \label{eq:policy-req}
\end{align}
Capacity and QoS constraints are
\begin{align}
  &\sum_{i=1}^{N}x_{ij}b_i\leq \min(cap_j,bh_j), &&\forall j,
  \label{eq:capacity}\\
  &b_i^{min}\leq b_i\leq b_i^{max},\quad
  \ell_i\leq\ell_i^{max},\quad
  \delta_i\leq\delta_i^{max}.
  \label{eq:qos}
\end{align}
Edge feasibility follows Eq.~\eqref{eq:edge-feasibility}.

\section{Interference Graph and Edge-Aware Attention}
\label{sec:interference-graph}

The graph state is
\begin{equation}
  G_t=(V_t,E_t,X_t,E_t^{feat}),
  \label{eq:edge-feature-graph}
\end{equation}
where $X_t$ stores node features and $E_t^{feat}$ stores edge features, including
interference weight, distance proxy, channel overlap, and utilisation coupling. The
AP-AP interference edge weight is
\begin{equation}
  w_{jl}(t)
  =
  \beta_1\mathbbm{1}[c_j=c_l]
  +\beta_2\frac{1}{(d_{jl}+\epsilon_d)^{\hat{\gamma}}}
  +\beta_3(\eta_j+\eta_l)
  +\beta_4O_{jl},
  \label{eq:edge-weight}
\end{equation}
with normalised terms, $\beta_r\geq0$, and $\sum_r\beta_r=1$.

Each scalar edge feature is min--max normalised within the scenario class before
training:
\begin{equation}
  \bar{x}_{r,jl}
  =
  \frac{x_{r,jl}-\min_{(u,v)\in E}x_{r,uv}}
  {\max_{(u,v)\in E}x_{r,uv}-\min_{(u,v)\in E}x_{r,uv}+\epsilon}.
  \label{eq:edge-feature-normalisation}
\end{equation}
The calibrated or learned coefficients $\beta_r$ are therefore comparable only within
a declared scenario class. They are not claimed to be globally transferable without
re-normalisation or recalibration.

The attention mechanism is edge-aware:
\begin{equation}
  \alpha_{uv}^{(l)}
  =
  \softmax_{u}
  \left(
  \operatorname{LeakyReLU}
  \left(
  \mathbf{a}^{(l)\top}
  \left[
  W^{(l)}h_v^{(l)}
  \Vert
  W^{(l)}h_u^{(l)}
  \Vert
  e_{uv}^{(l)}
  \right]
  \right)
  \right),
  \label{eq:edge-aware-attention}
\end{equation}
where $e_{uv}^{(l)}$ is the edge-feature vector. This corrects the earlier ambiguity
of hiding interference weights inside node features.

\section{Physics-Informed Path-Loss Scope}
\label{sec:pinn}

Indoor Wi-Fi is modelled with a log-distance multi-wall model:
\begin{equation}
  PL(d)
  =
  PL_0
  +10\hat{\gamma}\log_{10}\left(\frac{d}{d_0}\right)
  +\sum_{m=1}^{M_w}N_{w,m}L_{w,m}.
  \label{eq:multiwall}
\end{equation}
This model is validated against floor-plan/RSSI heatmap data, not terrain grid data.
The empirical acceptance target is $RMSE_{RSSI}\leq4$--$6\,\mathrm{dB}$ unless
stronger measurement evidence justifies a tighter bound.

For LEO or long-range rural backhaul, the physics-informed loss may use ITU-R
P.838/P.618 rain and earth-space propagation constraints~\cite{itu-p838,itu-p618}:
\begin{equation}
  \mathcal{L}_{PINN}
  =
  \mathcal{L}_{data}
  +\lambda_{phys}\mathcal{L}_{phys},
  \qquad
  \mathcal{L}_{phys}
  =
  \frac{1}{Q}
  \sum_{q=1}^{Q}
  (\hat{y}_q-y_{phys,q})^2.
  \label{eq:pinn-loss}
\end{equation}
Applying ITU-R rain attenuation to ordinary indoor Wi-Fi would be outside this
proposal's valid scope.

\section{Consistent Cost-Minimisation CMDP}
\label{sec:cmdp}

The slow controller is formulated as a cost-minimising CMDP:
\begin{equation}
  \mathcal{M}
  =
  (\mathcal{X},\Usafe(\cdot),P,C,\gamma),
  \qquad
  \gamma\in(0,1).
  \label{eq:cmdp}
\end{equation}
The state $x$ is a high-dimensional graph-embedding and telemetry summary, so the
Bellman equation is written with an expectation rather than a discrete summation:
\begin{equation}
  V^*(x)
  =
  \min_{u\in\Usafe(x)}
  \left[
  C(x,u)
  +\gamma
  \mathbb{E}_{x'\sim P(\cdot|x,u)}
  V^*(x')
  \right].
  \label{eq:bellman-corrected}
\end{equation}

If $\Usafe(x)=\emptyset$, the minimisation in Eq.~\eqref{eq:bellman-corrected}
is not evaluated over an empty action set. Instead, the controller invokes the
deterministic survivability fallback associated with the AP state $\omega_j$. In the
Online state, new admissions are rejected while existing valid sessions are preserved
only if they remain feasible. In the Degraded state, only previously authorised
subscribers with valid cached tokens may receive capped service and signed offline
accounting. In the Isolated state, no subscriber privilege escalation is permitted; the
AP either denies subscriber service or applies an explicitly configured guest-degraded
local policy.

The instantaneous cost is defined as
\begin{equation}
  C(x,u)
  =
  \omega_1\mathcal{I}(G)
  +\omega_2V_{SLA}
  +\omega_3O_{tun}
  +\omega_4O_{edge}
  +\omega_5O_{LEO}
  -\omega_6T_{norm}
  -\omega_7J(\tilde{\mathbf{b}}).
  \label{eq:cost-function}
\end{equation}
Here $T_{norm}$ is normalised delivered throughput. Latency, loss, and minimum-rate
violations are not separately embedded inside $T_{norm}$ because they are already
represented through $V_{SLA}$. This prevents the cost function from double-counting
the same service degradation through both utility and SLA-violation terms.

\begin{table}[H]
\centering
\caption{Operational definitions of CMDP cost terms.}
\label{tab:cmdp-cost-terms}
\begin{tabularx}{\textwidth}{p{0.16\linewidth}X p{0.24\linewidth}}
\toprule
\textbf{Term} & \textbf{Operational definition} & \textbf{Evidence source} \\
\midrule
$\mathcal{I}(G)$ & Normalised graph-interference cost computed from active AP--AP interference edge weights. & Radio-topology matrix. \\
$V_{SLA}$ & Fraction of subscriber-epochs violating latency, loss, or minimum-rate SLA bounds. & Slice-performance matrix. \\
$O_{tun}$ & Tunnel overhead ratio $(Bytes_{tunnel}-Bytes_{payload})/(Bytes_{payload}+\epsilon)$. & Slice-performance and gateway counters. \\
$O_{edge}$ & Normalised edge-resource pressure over CPU, RAM, IRQ, crypto, queueing, and telemetry load. & Edge-resource matrix. \\
$O_{LEO}$ & Normalised backhaul penalty for LEO path RTT, jitter, outage, handover, and power cost. & Backhaul-resilience matrix. \\
$T_{norm}$ & Delivered throughput normalised by subscriber entitlement or configured maximum. & Slice-performance matrix. \\
\bottomrule
\end{tabularx}
\end{table}

If a PPO-style implementation requires a reward signal, it is defined only as
\begin{equation}
  R(x,u)=-C(x,u).
  \label{eq:reward-negative-cost}
\end{equation}
This removes the earlier reward/cost sign-convention mismatch.

Because the state is represented by continuous telemetry and graph embeddings, any
PPO-style implementation uses function approximation. Therefore, tabular CMDP
convergence guarantees do not directly apply. Safety is enforced by the
state-dependent safe action filter rather than learned solely from reward penalties.

Fairness is measured using Jain's index over normalised subscriber entitlement:
\begin{equation}
  J(\tilde{\mathbf{b}})
  =
  \frac{
  \left(\sum_{i=1}^{N}\tilde{b}_i\right)^2
  }
  {
  N\sum_{i=1}^{N}\tilde{b}_i^2
  },
  \qquad
  \tilde{b}_i=\frac{b_i}{b_i^{max}}.
  \label{eq:jain}
\end{equation}

\section{LEO Backhaul Resilience Cost}
\label{sec:leo-model}

For backhaul mode $h$, effective capacity is
\begin{equation}
  bh_{j,t}^{eff}=A_{j,t}^{h}C_{j,t}^{h},
  \label{eq:bh-eff}
\end{equation}
where $A_{j,t}^{h}\in\{0,1\}$ is path availability and $C_{j,t}^{h}$ is usable
capacity.

For a LEO path, the cost uses normalised dimensionless terms:
\begin{equation}
  C_{LEO,j,t}^{cost}
  =
  \psi_1\frac{RTT_{j,t}}{RTT_{ref}}
  +\psi_2\frac{Jit_{j,t}}{Jit_{ref}}
  +\psi_3P_{out,j,t}
  +\psi_4Handover_{j,t}
  +\psi_5EnergyCost_{j,t}.
  \label{eq:leo-cost}
\end{equation}
The energy-cost term is defined as
\begin{equation}
  EnergyCost_{j,t}
  =
  1-\frac{PwrAvail_{j,t}}{PwrMax_j+\epsilon}.
  \label{eq:energy-cost}
\end{equation}
Thus, lower available terminal power increases the LEO cost. This prevents the
backhaul selector from accidentally treating high battery or power availability as a
penalty.

The selector chooses
\begin{equation}
  h^*_{j,t}
  =
  \argmin_{h\in\mathcal{H}}
  C_{h,j,t}^{cost}
  \quad
  \text{subject to }
  \pi\in\Pisafe.
  \label{eq:bh-selector}
\end{equation}